% \clearpage

% -------------------------
% Prompt 1: Segment Caption Integrator
% -------------------------
\begin{figure*}[t]
\centering
\begin{tcolorbox}[
  colback=blue!5!white,
  colframe=blue!75!black,
  width=\linewidth,
  arc=1mm,
  boxrule=0.8pt,
  title=\textbf{Segment Caption Integrator},
  colbacktitle=blue!75!black,
  coltitle=white,
  fonttitle=\bfseries\sffamily,
  halign title=center,
  top=4mm,
  bottom=4mm,
  left=3mm,
  right=3mm
]
\begin{lstlisting}[
  frame=none,
  breaklines=true,
  basicstyle=\ttfamily\small,
  columns=fullflexible,
  xleftmargin=0pt,
  xrightmargin=0pt,
  aboveskip=0pt,
  belowskip=0pt,
  gobble=0,
  breakindent=0pt,
  breakautoindent=false
]
You consolidate consecutive caption segments into one caption with no information loss.

Input has {segment_count} captions labeled S1 through S{segment_count}.
Return only valid JSON and do not output any other text.

Write in present tense with a neutral descriptive tone and keep the event order consistent with the input.

Use many short sentences and keep each sentence under 15 words.
Make each sentence atomic and keep one concrete observation per sentence.
Do not merge different events into one sentence.

Do not summarize and do not invent details.
Do not include timestamps or durations and avoid transition words such as then or afterward.

visual_caption describes only what is visible, including people, actions, objects, motion, background, lighting, camera viewpoint, and on screen text.
audio_caption describes only what is audible, including speech content, speaker tone, music, ambient noise, sound effects, silence, laughter, and applause.

{
  "visual_caption": "<short sentences describing only visible content>",
  "audio_caption": "<short sentences describing only audible content>"
}
\end{lstlisting}
\end{tcolorbox}
\caption{Prompt template for segment caption integration.}
\label{fig:caption_generation_prompt}
\end{figure*}


% -------------------------
% Prompt 2: Modality Specific Captioning
% -------------------------
\begin{figure*}[t]
\centering
\begin{tcolorbox}[
  colback=blue!5!white,
  colframe=blue!75!black,
  width=\linewidth,
  arc=1mm,
  boxrule=0.8pt,
  title=\textbf{Modality Specific Captioning},
  colbacktitle=blue!75!black,
  coltitle=white,
  fonttitle=\bfseries\sffamily,
  halign title=center,
  top=4mm,
  bottom=4mm,
  left=3mm,
  right=3mm
]
\begin{lstlisting}[
  frame=none,
  breaklines=true,
  basicstyle=\ttfamily\small,
  columns=fullflexible,
  xleftmargin=0pt,
  xrightmargin=0pt,
  aboveskip=0pt,
  belowskip=0pt,
  gobble=0,
  breakindent=0pt,
  breakautoindent=false
]
Analyze this {segment_duration}-second video segment and write exactly two parts as two lines.

VISUAL: Describe only what is visible, including people, objects, actions, scene context, colors, and motion.
AUDIO: Describe only what is audible, including speech content, speaker tone, music, sound effects, ambient sounds, and silence.

Keep both lines concise but informative and do not add any other text.
\end{lstlisting}
\end{tcolorbox}
\caption{Prompt template for modality specific captioning.}
\label{fig:modality_captioning_prompt}
\end{figure*}


% -------------------------
% Prompt 3: Moment Search Query Generation
% -------------------------
\begin{figure*}[t]
\centering
\begin{tcolorbox}[
  colback=blue!5!white,
  colframe=blue!75!black,
  width=\linewidth,
  arc=1mm,
  boxrule=0.8pt,
  title=\textbf{Moment Search Query Generation},
  colbacktitle=blue!75!black,
  coltitle=white,
  fonttitle=\bfseries\sffamily,
  halign title=center,
  top=4mm,
  bottom=4mm,
  left=3mm,
  right=3mm
]
\begin{lstlisting}[
  frame=none,
  breaklines=true,
  basicstyle=\ttfamily\small,
  columns=fullflexible,
  xleftmargin=0pt,
  xrightmargin=0pt,
  aboveskip=0pt,
  belowskip=0pt,
  gobble=0,
  breakindent=0pt,
  breakautoindent=false
]
Generate one short and realistic query that a viewer would type to revisit the exact moment behind the comment.

First infer whether the intent is visual or audio.
Extract comment_keywords as words or short phrases taken from the comment.
Select focus_keywords as the smallest subset that carries the retrieval intent.

Write one query that centers on focus_keywords and add one or two small cues that help localize the moment.
Use a conversational search style and write one sentence in 12 to 25 words.

For visual queries emphasize visible entities, actions, and scene details.
For audio queries include a short quote or a natural paraphrase and add sound cues when helpful.

Avoid timestamps and durations and avoid narration that describes watching the video.
Return valid JSON only and follow this schema.

{
  "modality": "visual" or "audio",
  "comment_keywords": ["phrases taken from the comment"],
  "focus_keywords": ["subset used to build the query"],
  "intent": "one sentence describing what the user wants to revisit",
  "query": "one natural query in 12 to 25 words",
  "rationale": "why the focus keywords and cues match the comment intent"
}

Example
Comment: "Love this quote, Jillian isn't sick she's a dancer!"
{
  "modality": "audio",
  "comment_keywords": ["Love this quote", "Jillian isn't sick", "dancer"],
  "focus_keywords": ["Jillian isn't sick", "dancer"],
  "intent": "Find the moment when the memorable Jillian line is spoken",
  "query": "the moment someone says Jillian isn't sick she's a dancer and the audience reacts",
  "rationale": "The key phrase comes from the comment and the reaction cue helps localization"
}
\end{lstlisting}
\end{tcolorbox}
\caption{Prompt template for moment search query generation.}
\label{fig:search_query_generation_prompt}
\end{figure*}

\begin{tcolorbox}[
  colback=blue!5!white,
  colframe=blue!75!black,
  width=\textwidth,
  arc=1mm,
  boxrule=0.8pt,
  boxsep=1mm,
  title=\textbf{Modality Specific Captioning},
  colbacktitle=blue!75!black,
  coltitle=white,
  fonttitle=\bfseries\sffamily,
  halign title=center,
  top=2mm,
  bottom=2mm,
  left=2mm,
  right=2mm
]
\small
\setlength{\parindent}{0pt}
\begin{lstlisting}[
  frame=none,
  breaklines=true,
  basicstyle=\ttfamily\small,
  columns=fullflexible,
  aboveskip=0pt,
  belowskip=0pt
]
Analyze this {segment_duration}-second video segment and write exactly two parts as two lines.

VISUAL: Describe only what is visible, including people, objects, actions, scene context, colors, and motion.
AUDIO: Describe only what is audible, including speech content, speaker tone, music, sound effects, ambient sounds, and silence.

Keep both lines concise but informative and do not add any other text.
\end{lstlisting}
\end{tcolorbox}